\documentclass[]{report}

\usepackage{import}

\import{../}{settings.tex}

\title{Notions}
\author{Notions ou propriété en plus dans le cadre de l'aggrégation}
\date{}

\begin{document}

    \maketitle
    \tableofcontents

    \chapter{Algèbre}

    \begin{prop}{Continuité d'un morphisme de groupe}{}
        Soit $\funto{f}{G}{H}$ un morphisme de groupes. Si $f$ est continue en $a\in G$, alors $f$ est continue en tout point de $G$.\\
        Conséquence: si $f$ est continue en $e_G$, alors $f$ est continue sur $G$.
    \end{prop}

    \begin{proof}{Continuité d'un morphisme de groupe}{}
        Supposons que $\funto{f}{G}{H}$ est un morphisme de groupes qui est continue en $a\in G$.
        \[\forall W\in \Vr(f(a)),\exists V\in \Vr(a), f(V)\subset W\]
        Supposons d'abord que $a=e_G$. Soit $y\in G$ et $W\in \Vr(f(y))$. En posant,
        \[W'=f(y)^{-1}\cdot W\]
        On a que $f(y)^{-1}\cdot f(y)=e_H\in W'$ et donc que $W'$ est un voisinage de $e_H$, il existe donc $U'\in \Vr(e_G)$, tel que $f(U')\subset W'$. On translate par $y$ en posant
        \[U=y\cdot U'\]
        qui est un voisinage de $y\cdot e_G=y$. On vérifie alors que $f(U)\subset W$. Soit $x\in U, x=yu$,
        \[f(x)=f(y)f(u)\in f(y)W'=f(y)f(y)^{-1}W=W\]
        Donc Si $f$ est continue en $e_G$, alors $f$ est continue sur $G$.\\
        Supposons maintenant que $a\neq 0$, et montrons que $f$ est continue en $e_G$. Soit $W$ voisinage de $e_H$, on pose
        \[W'=W\cdot f(a)\]
        comme $W$ est un voisinage de $e_H$, alors $f(a)\in W'$, donc il existe $V\in \Vr(a)$ tel que
        \[f(V)\subset W'\]
        On pose $V'=V\cdot a^{-1}$ est un voisinage de $a a^{-1}=e_G$. On vérifie que $f(V')\subset W$. Soit $x\in V'$, donc $xa\in V$,
        \[f(x)=f(xa) f(a)^{-1}\in W' \cdot f(a)^{-1}=(W\cdot f(a))\cdot f(a)^{-1}=W\]
        On a donc que $f(V)\subset W$ et donc que $f$ est continue en $e_G$ et donc sur $G$
        Il reste simplement à montrer que les translations sont continues, ce qui dépend de la topologie appliquée sur $G$ et $H$
    \end{proof}

    \section{Algèbre Linéaire}

    \section{Algèbre Générale}

    \chapter{Analyse}

    \chapter{Probabilité}
\end{document}